% Mustervorlage für Abschlussarbeiten am LS DBIS
% Author: Alexander Stahl (stahlale@b-tu.de)
% Zuletzt aktualisiert am 06.06.2023

% Hinweis: Die Daten zur Arbeit speichern Sie in der Datei "data.tex".

% Diese Datei enthält Einstellungen und Paketimporte.
% Zur besseren Ordnung sollten Sie weitere Einstellungen bzw. Importe dieser Datei anfügen.

% Paketimporte
\documentclass[11pt,a4paper,german]{scrreprt}
\usepackage[T1]{fontenc}  % Encoding
\usepackage[utf8]{inputenc}  % Encoding
\usepackage{babel}  % Unterstützung verschiedener Sprachen
\usepackage{graphicx}  % Einbindung von Grafiken
\usepackage{amsmath}  % Verschiedene Symbole
\usepackage{amssymb}  % Verschiedene Symbole
\usepackage{latexsym}  % Verschiedene Symbole
\usepackage{booktabs}  % Schönere Tabellen
\usepackage{multicol}  % Flexiblere Tabellen
\usepackage[printonlyused]{acronym}  % Abkürzungsverwaltung
\usepackage{hyperref}  % Klickbare URLs, Titel und PDF-Metadaten
\usepackage{lipsum}  % Platzhaltertext (kann entfernt werden)
\usepackage[backend=bibtex,style=ieee]{biblatex}  % Moderne Literaturverwaltung

% Einstellungen
\hypersetup{
    colorlinks,
    citecolor=black,
    filecolor=black,
    linkcolor=black,
    urlcolor=black
}

\makeatletter

\newcommand\frontmatter{
	\cleardoublepage
	\pagenumbering{roman}}

\newcommand\mainmatter{
	\cleardoublepage
	\pagenumbering{arabic}}

\newcommand\backmatter{
  \if@openright
    \cleardoublepage
  \else
    \clearpage
  \fi
   }

\makeatother

% Laden der Literaturdaten
\addbibresource{quellen.bib}

% Definition der Dokumentendaten
% Zur besseren Kapselung wurde die Eingabe Ihrer Daten in diese Datei ausgelagert.
% Bitte einfach ausfüllen.

\newcommand{\TitelDeutsch}{Titel der Arbeit} % Hier zentral angeben
\newcommand{\TitelEnglisch}{Title of the Thesis} % Übersetzten Titel auch angeben
\newcommand{\ArbeitsTyp}{Bachelorarbeit/Masterarbeit} % unzutreffendes hier löschen
\newcommand{\AuthorName}{Max Mustermann} % Hier ändern
\newcommand{\MatrikelNr}{(hier einfügen)}
\newcommand{\Studiengang}{(hier einfügen)}
\newcommand{\DatumThemenausgabe}{(hier einfügen)}
\newcommand{\DatumAbgabe}{(hier einfügen)}
\newcommand{\BetreuerEins}{Prof. Dr.-Ing. Ingo Schmitt}
\newcommand{\BetreuerZwei}{(ggf. einfügen)} % Bitte immer akademische Titel angeben
  % In dieser Datei die Daten der Arbeit eingeben.

\titlehead{{Brandenburgische Technische Universität Cottbus-Senftenberg\\ Institut für Informatik\\ Fachgebiet Datenbank- und Informationssysteme}}
\subject{Bachelorarbeit/Masterarbeit\\ \vspace{0.5cm}\includegraphics[width=0.5\textwidth]{img/btu-logo.pdf}}
\title{\TitelDeutsch \\ {\normalsize \TitelEnglisch}}
\author{\AuthorName \\ MatrikelNr.: \MatrikelNr \\Studiengang: \Studiengang}
\date{\textit{Datum der Themenausgabe: \DatumThemenausgabe} \\ \textit{Datum der Abgabe: \DatumAbgabe}}
\publishers{1. Betreuer: \BetreuerEins \\ 2. Betreuer: \BetreuerZwei}

\hypersetup{pdfauthor={\AuthorName}}
\hypersetup{pdftitle={\TitelDeutsch}}
  % Header beinhaltet alle Einstellungen und Pakete. Bei Bedarf erweitern.

\begin{document}
  \frontmatter  % Kennzeichnet den vorderen Teil der Arbeit
  \maketitle  % Titelseite generieren

  % Eidesstattliche Erklärung über Selbstständigkeit
  \section*{Eidesstattliche Erklärung}
Der Verfasser der Arbeit erklärt hiermit Folgendes:

\begin{enumerate}
  \item Die vorliegende Arbeit wurde selbständig, ohne fremde Hilfe und ohne Benutzung anderer als der angegebenen Hilfsmittel angefertigt. Es wurden keine automatisierten Schreibhilfen oder künstliche Intelligenz zur Erstellung dieses Werkes verwendet.
  \item Die aus fremden Quellen (einschließlich elektronischer Quellen) direkt oder indirekt übernommenen Informationen und Gedanken sind ausnahmslos als solche kenntlich gemacht. Wörtlich und inhaltlich verwendete Quellen wurden entsprechend den anerkannten Regeln wissenschaftlichen Arbeitens zitiert.
  \item Die Arbeit ist nicht in dieser oder vergleichbarer Form (auch nicht auszugsweise) im Rahmen einer anderen Prüfung bei einer anderen Hochschule vorgelegt oder publiziert worden.
  \item Der Verfasser ist damit einverstanden, dass die Arbeit mit Hilfe eines Plagiatserkennungsdienstes auf enthaltene Plagiate überprüft wird.
\end{enumerate}

Dem Verfasser ist bewusst, dass falsche Angaben und Plagiate schwerwiegende Konsequenzen nach sich ziehen können, einschließlich einer Aberkennung des akademischen Grades. Er ist sich der Verantwortung bewusst, die mit dem Verfassen einer wissenschaftlichen Arbeit einhergeht, und hat die Arbeit nach bestem Wissen und Gewissen verfasst.


\vspace{2cm}

\noindent
\begin{tabular}{lcl}
  .................................... & \hspace{.4\textwidth} & ....................................\\
  Ort, Datum &  & Unterschrift

\end{tabular}

  \newpage

  % Danksagung (Nur falls gewünscht. Ggf. Kommentarzeichen entfernen.)
  % \section*{Danksagung}
Hier ggf. Personen danken, die indirekt bei der Arbeit unterstützt haben, z. B. durch finanzielles Unterstützen, Mutmachen, Korrekturlesen, etc.
  % \newpage

  % Abstract bzw. Zusammenfassung der Arbeit
  \section*{Abstract}
Verfassen Sie hier Ihren Abstract. Dabei handelt es sich um eine prägnante Zusammenfassung Ihrer Abschlussarbeit, die Ihren Lesern einen Überblick über den Inhalt, die Ziele, die Methodik, die Ergebnisse und die Schlussfolgerungen der Arbeit bietet. Vermeiden Sie es, Informationen zu geben, die in der Arbeit selbst nicht auftauchen und halten Sie sich an eine Länge von ca. 100-300 Wörtern. Der Abstract muss aus einem einzigen Absatz bestehen.
  \newpage

  % Automatische Generierung verschiedener Verzeichnisse
  % Inhaltsverzeichnis
  \tableofcontents

  % Abbildungsverzeichnis
  \cleardoublepage\phantomsection
  \addcontentsline{toc}{chapter}{\listfigurename}\listoffigures

  % Tabellenverzeichnis:
  \cleardoublepage\phantomsection
  \addcontentsline{toc}{chapter}{\listtablename}\listoftables

  % Abkürzungsverzeichnis:
  \cleardoublepage\phantomsection
  \addcontentsline{toc}{chapter}{Abkürzungsverzeichnis}
  \chapter*{Abkürzungsverzeichnis}
  \begin{acronym}[Bash]
	% Hier wie folgt Abkürzungen definieren: \acro{Abkürzung}{Volles Wort}
	% Anschließend kann die Abkürzung im Text so benutzt werden: \ac{Abkürzung}
	% Dadurch wird die Abkürzung immer an der ersten Nennung in Klammern definiert.

	%Beispiel
	\acro{VM}{Virtuelle Maschine}
	\acrodefplural{VM}[VMs]{Virtuelle Maschinen}
\end{acronym}


  \mainmatter  % Kennzeichnet den Hauptteil der Arbeit

  % Es folgt der Haupttext der Arbeit. Passen Sie die Überschriften der Kapitel entsprechend an.
  % Beachten Sie, dass die hier angegebene Gliederung nur ein Vorschlag ist.
  % Außerdem wird empfohlen, jedes Kapitel in eine eigene Datei auszulagern und hier nur mittels \input einzubinden.

  \chapter{Einleitung}
  % Beispiel: \input{kapitel1.tex}
  \lipsum[1-4]  % Platzhaltertext (bitte entfernen)
  \chapter{Übersicht über das Themengebiet}
  \lipsum[1-4]  % Platzhaltertext (bitte entfernen)
  \chapter{Eigener Ansatz}
  \lipsum[1-5]  % Platzhaltertext (bitte entfernen)
  \chapter{Implementierung des Ansatzes}
  \lipsum[1-4]  % Platzhaltertext (bitte entfernen)
  \chapter{Experimentelle bzw. Formale Evaluierung}
  \lipsum[1-4]  % Platzhaltertext (bitte entfernen)
  \chapter{Zusammenfassung und Ausblick}
  \lipsum[1-3]  % Platzhaltertext (bitte entfernen)

  \backmatter  % Kennzeichnet den hinteren Teil der Arbeit
  \appendix  % Anhang

  \cleardoublepage
  \addcontentsline{toc}{chapter}{Literaturverzeichnis}
  % \nocite{*}  % Hier das Kommentarzeichen entfernen, um alle (auch ungenutzte) Quellen auszugeben.
  \printbibliography[title={Literaturverzeichnis}]  % Generiert das Literaturverzeichnis

  \vfill
  \begin{center}
   \emph{Onlineressourcen wurden im <Monat> <Jahr> auf ihre Verfügbarkeit hin überprüft.}
  \end{center}

\end{document}
